\section{Título y Objetivos}
\label{sec:titulo/objetivos}

\noindent \textbf{Práctica 1:} Verificación experimental de la Ley de Ohm mediante la medición de corriente y voltaje en una resistencia metálica, y determinación de su resistencia por ajuste lineal.

\noindent \textbf{Práctica 2:} Medición de resistencias mediante un multímetro digital (modo ohmímetro), identificación del valor nominal por código de colores, y verificación de compatibilidad con la tolerancia declarada por el fabricante.

\noindent \textbf{Práctica 3:} Análisis de la influencia del aparato de medida sobre el circuito, verificación de la Segunda Ley de Kirchhoff en un divisor resistivo y evaluación del efecto de la resistencia de entrada del voltímetro.

\noindent Los objetivos específicos de las tres prácticas son:
\begin{itemize}
    \item Realizar mediciones directas de corriente, voltaje y resistencia con instrumentos de laboratorio.
    \item Determinar la resistencia de un conductor por regresión lineal sobre datos de corriente y voltaje.
    \item Estimar incertezas instrumentales a partir de las especificaciones del fabricante y propagarlas correctamente.
    \item Verificar leyes de circuitos (Ley de Ohm, Segunda Ley de Kirchhoff) contrastando modelo teórico con datos experimentales.
    \item Expresar todos los resultados con sentido físico (valor representativo $\pm$ incerteza absoluta, con unidades).
\end{itemize}
